\section{Combination Sum II}
\label{sec:rec-combination-sum-2}

\problemstatement{Given an array $\textit{candidates}$ that \textbf{may
contain duplicate values} and a target, return every unique combination
that sums to the target, where \textbf{each individual element of the
array may be used at most once} (no reuse), and the output must contain
no duplicate combinations.}

\begin{patternbox}
This problem directly combines two techniques already derived in this
chapter: Section~\ref{sec:rec-subset-sum-2}'s ``sort, then skip a choice
equal to the previous choice at the same level'' duplicate-avoidance, and
Section~\ref{sec:rec-combination-sum}'s general combination-building
recursion -- but \textbf{without} that section's ``stay at the same
index'' reuse rule, since each element here may only be used once.
\end{patternbox}

\subsection*{Understanding the problem carefully}

Because reuse is forbidden, the recursion always advances past the
current index after picking it (exactly like
Section~\ref{sec:rec-subset-sum-2}, never like
Section~\ref{sec:rec-combination-sum}). Because duplicate values are
present, we sort first and skip same-level duplicate choices, exactly as
in Section~\ref{sec:rec-subset-sum-2}, to avoid generating the same
combination more than once. One further optimization becomes available
specifically because we additionally have a numeric target and a sorted
array: the moment a candidate's value alone exceeds the remaining
target, every later candidate (being at least as large, since the array
is sorted) must also exceed it, so the entire rest of the loop at this
level can be abandoned immediately, rather than checked one by one.

\begin{ideabox}[Sort, Skip Same-Level Duplicates, and Break Early]
Sort $\textit{candidates}$. $\textsc{Solve}(\textit{start},
\textit{remaining}, \textit{current})$: if $\textit{remaining}=0$:
record $\textit{current}$; return. For $i$ from $\textit{start}$ to
$n-1$: if $i>\textit{start}$ and $\textit{candidates}[i] =
\textit{candidates}[i-1]$: skip (same-level duplicate). If
$\textit{candidates}[i] > \textit{remaining}$: \textbf{break} out of the
loop entirely (pruning -- no further $i$ in this sorted loop can help).
Otherwise: append $\textit{candidates}[i]$; recurse
$\textsc{Solve}(i+1, \textit{remaining} - \textit{candidates}[i],
\textit{current})$; remove (undo).
\end{ideabox}

\subsection*{Why breaking (not just skipping) on an over-large candidate is valid}

This is a direct consequence of sortedness, not a heuristic: since
$\textit{candidates}$ is sorted in increasing order, once
$\textit{candidates}[i] > \textit{remaining}$, every subsequent
$\textit{candidates}[j]$ for $j>i$ satisfies $\textit{candidates}[j] \ge
\textit{candidates}[i] > \textit{remaining}$ too -- so none of them could
possibly lead to a valid combination either, at this level. Breaking the
loop entirely (rather than merely skipping this one $i$ and checking the
next) is therefore a safe pruning step that discards an entire suffix of
guaranteed-failing branches in $O(1)$, rather than wastefully visiting
and individually rejecting each one.

\subsection*{Algorithm}

\begin{enumerate}
  \item Sort $\textit{candidates}$.
  \item $\textsc{Solve}(\textit{start}, \textit{remaining}, \textit{current})$:
  \begin{enumerate}
    \item If $\textit{remaining}=0$: record $\textit{current}$; return.
    \item For $i$ from $\textit{start}$ to $n-1$:
    \begin{itemize}
      \item If $i>\textit{start}$ and $\textit{candidates}[i] =
            \textit{candidates}[i-1]$: continue (skip).
      \item If $\textit{candidates}[i] > \textit{remaining}$: break.
      \item Append $\textit{candidates}[i]$; recurse
            $\textsc{Solve}(i+1, \textit{remaining} -
            \textit{candidates}[i], \textit{current})$; remove (undo).
    \end{itemize}
  \end{enumerate}
  \item Call $\textsc{Solve}(0, \textit{target}, [\,])$.
\end{enumerate}

\subsection*{Dry run}

\dryrun{Take $\textit{candidates} = [1,1,2]$ (sorted),
$\textit{target}=3$.}

\begin{itemize}
  \item $\textsc{Solve}(0,3,[])$: $i=0$ ($1$): append, recurse
        $\textsc{Solve}(1,2,[1])$. Undo.
  \begin{itemize}
    \item $\textsc{Solve}(1,2,[1])$: $i=1$ ($1$): $i=\textit{start}$,
          no skip check needed. Append, recurse
          $\textsc{Solve}(2,1,[1,1])$. Undo.
    \begin{itemize}
      \item $\textsc{Solve}(2,1,[1,1])$: $i=2$ ($2$): $2 > 1$,
            \textbf{break} immediately.
    \end{itemize}
    \item $i=2$ ($2$) at this level: append, recurse
          $\textsc{Solve}(3,0,[1,2])$: $\textit{remaining}=0$! Record
          $[1,2]$.
  \end{itemize}
  \item Back at top level: $i=1$ ($1$): $i>\textit{start}=0$ and
        $\textit{candidates}[1]=\textit{candidates}[0]$: skip.
  \item $i=2$ ($2$): append, recurse $\textsc{Solve}(3,1,[2])$:
        $\textit{remaining}=1 \ne 0$, and the loop from $i=3$ to $n-1$
        is empty (no candidates left); this branch records nothing.
\end{itemize}

Recorded combinations: $[1,1,1]$? -- wait, recheck: we recorded $[1,2]$
once. The full correct output for this classic instance is
$\{[1,1,1]\text{-like attempts that failed}, [1,2]\}$; tracing carefully
as above, the only successful combination found is $[1,2]$, matching the
expected unique result for target $3$ from $\{1,1,2\}$ (the combination
$1+1+2=4 \ne 3$ does not qualify; $1+2=3$ does, exactly once despite two
copies of $1$ being available).

\subsection*{C++ implementation}

\begin{lstlisting}
#include <bits/stdc++.h>
using namespace std;

void solve(int start, int remaining, vector<int>& candidates,
           vector<int>& current, vector<vector<int>>& result) {
    if (remaining == 0) {
        result.push_back(current);
        return;
    }
    for (int i = start; i < (int)candidates.size(); i++) {
        if (i > start && candidates[i] == candidates[i - 1]) continue;
        if (candidates[i] > remaining) break;

        current.push_back(candidates[i]);
        solve(i + 1, remaining - candidates[i], candidates, current, result);
        current.pop_back();
    }
}

vector<vector<int>> combinationSum2(vector<int>& candidates, int target) {
    sort(candidates.begin(), candidates.end());
    vector<vector<int>> result;
    vector<int> current;
    solve(0, target, candidates, current, result);
    return result;
}

int main() {
    vector<int> candidates = {1, 1, 2};
    for (auto& comb : combinationSum2(candidates, 3)) {
        cout << "[ ";
        for (int x : comb) cout << x << " ";
        cout << "] ";
    }
    cout << endl;
    return 0;
}
\end{lstlisting}

\begin{complexitybox}
\textbf{Time Complexity.} $O(2^n)$ in the worst case (every subset
potentially checked), though sorting plus the early-break pruning
typically eliminates a substantial fraction of branches in practice; an
additional $O(n \log n)$ for the initial sort.
\textbf{Space Complexity.} $O(n)$ for the recursion depth.
\end{complexitybox}

\begin{takeawaybox}
Combination Sum II is a direct composition of two already-derived
techniques -- no-reuse advancing recursion plus sort-and-skip duplicate
avoidance -- with one additional sorted-array pruning trick (breaking
early once a candidate exceeds the remaining target) layered on top.
Recognizing a new problem as ``technique A combined with technique B''
is, again, the fastest path to a correct solution.
\end{takeawaybox}
